% Options for packages loaded elsewhere
\PassOptionsToPackage{unicode}{hyperref}
\PassOptionsToPackage{hyphens}{url}
\PassOptionsToPackage{dvipsnames,svgnames,x11names}{xcolor}
%
\documentclass[
  10pt,
]{article}
\usepackage{amsmath,amssymb}
\usepackage{iftex}
\ifPDFTeX
  \usepackage[T1]{fontenc}
  \usepackage[utf8]{inputenc}
  \usepackage{textcomp} % provide euro and other symbols
\else % if luatex or xetex
  \usepackage{unicode-math} % this also loads fontspec
  \defaultfontfeatures{Scale=MatchLowercase}
  \defaultfontfeatures[\rmfamily]{Ligatures=TeX,Scale=1}
\fi
\usepackage{lmodern}
\ifPDFTeX\else
  % xetex/luatex font selection
  \setmainfont[]{DejaVu Serif}
  \setsansfont[]{DejaVu Sans}
  \setmonofont[]{DejaVu Sans Mono}
  \setmathfont[]{Latin Modern Math}
\fi
% Use upquote if available, for straight quotes in verbatim environments
\IfFileExists{upquote.sty}{\usepackage{upquote}}{}
\IfFileExists{microtype.sty}{% use microtype if available
  \usepackage[]{microtype}
  \UseMicrotypeSet[protrusion]{basicmath} % disable protrusion for tt fonts
}{}
\makeatletter
\@ifundefined{KOMAClassName}{% if non-KOMA class
  \IfFileExists{parskip.sty}{%
    \usepackage{parskip}
  }{% else
    \setlength{\parindent}{0pt}
    \setlength{\parskip}{6pt plus 2pt minus 1pt}}
}{% if KOMA class
  \KOMAoptions{parskip=half}}
\makeatother
\usepackage{xcolor}
\usepackage[margin=2.0cm]{geometry}
\usepackage{color}
\usepackage{fancyvrb}
\newcommand{\VerbBar}{|}
\newcommand{\VERB}{\Verb[commandchars=\\\{\}]}
\DefineVerbatimEnvironment{Highlighting}{Verbatim}{commandchars=\\\{\}}
% Add ',fontsize=\small' for more characters per line
\newenvironment{Shaded}{}{}
\newcommand{\AlertTok}[1]{\textcolor[rgb]{1.00,0.00,0.00}{\textbf{#1}}}
\newcommand{\AnnotationTok}[1]{\textcolor[rgb]{0.38,0.63,0.69}{\textbf{\textit{#1}}}}
\newcommand{\AttributeTok}[1]{\textcolor[rgb]{0.49,0.56,0.16}{#1}}
\newcommand{\BaseNTok}[1]{\textcolor[rgb]{0.25,0.63,0.44}{#1}}
\newcommand{\BuiltInTok}[1]{\textcolor[rgb]{0.00,0.50,0.00}{#1}}
\newcommand{\CharTok}[1]{\textcolor[rgb]{0.25,0.44,0.63}{#1}}
\newcommand{\CommentTok}[1]{\textcolor[rgb]{0.38,0.63,0.69}{\textit{#1}}}
\newcommand{\CommentVarTok}[1]{\textcolor[rgb]{0.38,0.63,0.69}{\textbf{\textit{#1}}}}
\newcommand{\ConstantTok}[1]{\textcolor[rgb]{0.53,0.00,0.00}{#1}}
\newcommand{\ControlFlowTok}[1]{\textcolor[rgb]{0.00,0.44,0.13}{\textbf{#1}}}
\newcommand{\DataTypeTok}[1]{\textcolor[rgb]{0.56,0.13,0.00}{#1}}
\newcommand{\DecValTok}[1]{\textcolor[rgb]{0.25,0.63,0.44}{#1}}
\newcommand{\DocumentationTok}[1]{\textcolor[rgb]{0.73,0.13,0.13}{\textit{#1}}}
\newcommand{\ErrorTok}[1]{\textcolor[rgb]{1.00,0.00,0.00}{\textbf{#1}}}
\newcommand{\ExtensionTok}[1]{#1}
\newcommand{\FloatTok}[1]{\textcolor[rgb]{0.25,0.63,0.44}{#1}}
\newcommand{\FunctionTok}[1]{\textcolor[rgb]{0.02,0.16,0.49}{#1}}
\newcommand{\ImportTok}[1]{\textcolor[rgb]{0.00,0.50,0.00}{\textbf{#1}}}
\newcommand{\InformationTok}[1]{\textcolor[rgb]{0.38,0.63,0.69}{\textbf{\textit{#1}}}}
\newcommand{\KeywordTok}[1]{\textcolor[rgb]{0.00,0.44,0.13}{\textbf{#1}}}
\newcommand{\NormalTok}[1]{#1}
\newcommand{\OperatorTok}[1]{\textcolor[rgb]{0.40,0.40,0.40}{#1}}
\newcommand{\OtherTok}[1]{\textcolor[rgb]{0.00,0.44,0.13}{#1}}
\newcommand{\PreprocessorTok}[1]{\textcolor[rgb]{0.74,0.48,0.00}{#1}}
\newcommand{\RegionMarkerTok}[1]{#1}
\newcommand{\SpecialCharTok}[1]{\textcolor[rgb]{0.25,0.44,0.63}{#1}}
\newcommand{\SpecialStringTok}[1]{\textcolor[rgb]{0.73,0.40,0.53}{#1}}
\newcommand{\StringTok}[1]{\textcolor[rgb]{0.25,0.44,0.63}{#1}}
\newcommand{\VariableTok}[1]{\textcolor[rgb]{0.10,0.09,0.49}{#1}}
\newcommand{\VerbatimStringTok}[1]{\textcolor[rgb]{0.25,0.44,0.63}{#1}}
\newcommand{\WarningTok}[1]{\textcolor[rgb]{0.38,0.63,0.69}{\textbf{\textit{#1}}}}
\usepackage{longtable,booktabs,array}
\usepackage{calc} % for calculating minipage widths
% Correct order of tables after \paragraph or \subparagraph
\usepackage{etoolbox}
\makeatletter
\patchcmd\longtable{\par}{\if@noskipsec\mbox{}\fi\par}{}{}
\makeatother
% Allow footnotes in longtable head/foot
\IfFileExists{footnotehyper.sty}{\usepackage{footnotehyper}}{\usepackage{footnote}}
\makesavenoteenv{longtable}
\setlength{\emergencystretch}{3em} % prevent overfull lines
\providecommand{\tightlist}{%
  \setlength{\itemsep}{0pt}\setlength{\parskip}{0pt}}
\setcounter{secnumdepth}{-\maxdimen} % remove section numbering
% doc_header.tex -- LaTeX preamble additions used by build_docs.sh for every
% Markdown -> LaTeX -> PDF conversion (supplements, README, submission files).
\usepackage{etoolbox}
% longtable may not be loaded by pandoc when a document has no table
\usepackage{longtable}
% bold vectors: \bm -> \symbf under unicode-math (XeLaTeX), bm package otherwise
\makeatletter
\@ifpackageloaded{unicode-math}{\providecommand{\bm}[1]{\symbf{#1}}}{\usepackage{bm}}
\makeatother
% smaller type in long tables so that wide audit tables fit the text width;
% inside a table a line break is allowed after an underscore of plain text
% (pandoc writes it as \_), so that long identifiers wrap inside their column
% rows are set a little apart, so that symbols with indices do not touch
\AtBeginEnvironment{longtable}{\footnotesize\renewcommand{\arraystretch}{1.45}%
  \let\sptunderscore\_\def\_{\sptunderscore\allowbreak}}
\setlength{\LTcapwidth}{\textwidth}
% allow line breaks after slashes and hyphens inside table cells
\usepackage{xurl}
% code blocks: break over-long lines instead of running into the margin.
% Plain ``` blocks are typeset by pandoc as {verbatim}, highlighted ones
% (```bash) as {Highlighting}; both are re-declared as breakable fancyvrb
% environments.  The continuation line is marked with a small hooked arrow.
\usepackage{fvextra}
\RecustomVerbatimEnvironment{verbatim}{Verbatim}{breaklines,breakanywhere,fontsize=\small}
\makeatletter
\@ifundefined{Highlighting}{}{%
  \RecustomVerbatimEnvironment{Highlighting}{Verbatim}{breaklines,breakanywhere,fontsize=\small,commandchars=\\\{\}}}
\makeatother
% inline code (pandoc writes it as \texttt{...}): never hyphenate -- an inserted
% hyphen could be read as part of a file name -- but allow a line break after
% an underscore, slash, comma, colon or equals sign, so that long paths and
% column lists wrap inside the text width and inside table columns.
% Inside a formula (e.g. the name of a notebook symbol in an equation) the same
% command must first leave math mode, which \text of amsmath does.
\ExplSyntaxOn
\tl_new:N \l__spt_code_tl
\cs_new_protected:Npn \__spt_texttt:n #1
  {
    \tl_set:Nn \l__spt_code_tl {#1}
    \regex_replace_all:nnN { ([\/\,\:\=]) } { \1 \c{allowbreak} } \l__spt_code_tl
    \regex_replace_all:nnN { \c{_} } { \c{_} \c{allowbreak} } \l__spt_code_tl
    \group_begin: \ttfamily \hyphenchar \font = -1 \scan_stop: \l__spt_code_tl \group_end:
  }
\RenewDocumentCommand \texttt { m }
  {
    \mode_if_math:TF { \text { \__spt_texttt:n {#1} } } { \__spt_texttt:n {#1} }
  }
\ExplSyntaxOff
% A figure is written in the Markdown source as an image paragraph followed by
% a caption paragraph ("**Figure S1.** ..."), so that the caption is also
% visible on GitHub.  An image that starts a paragraph is centered.
\usepackage{graphicx}% pandoc loads it only for documents that contain an image
\NewCommandCopy\sptincludegraphics\includegraphics
\RenewDocumentCommand\includegraphics{O{} m}{%
  \ifvmode\begin{center}\sptincludegraphics[#1]{#2}\end{center}%
  \else\sptincludegraphics[#1]{#2}\fi}
% keep figures where they are referenced
\usepackage{float}
\floatplacement{figure}{H}
% running header with the document title
\usepackage{fancyhdr}
\pagestyle{fancy}
\fancyhf{}
\fancyhead[L]{\small PROJECT spacetimes, Revision\_Deflation --- documentation and provenance}
\fancyhead[R]{\small\thepage}
\renewcommand{\headrulewidth}{0.4pt}
% 16x16 matrices: tighter columns so that a full matrix fits the text width
\setlength{\arraycolsep}{3.2pt}
% The documents of this project contain long tagged symbols, file paths and
% identifiers that cannot be hyphenated.  A paragraph is therefore allowed to
% have looser lines before TeX gives up and lets a line stick out of the margin.
\AtBeginDocument{\tolerance=2000 \emergencystretch=5em \hbadness=2000}
\ifLuaTeX
  \usepackage{selnolig}  % disable illegal ligatures
\fi
\IfFileExists{bookmark.sty}{\usepackage{bookmark}}{\usepackage{hyperref}}
\IfFileExists{xurl.sty}{\usepackage{xurl}}{} % add URL line breaks if available
\urlstyle{same}
\hypersetup{
  pdftitle={RP8 --- The source code of Revision\_Deflation: what computes what, and how to compile and run it},
  pdfauthor={Prepared for Patrick Lee Nash, Ph.D.},
  colorlinks=true,
  linkcolor={blue},
  filecolor={Maroon},
  citecolor={blue},
  urlcolor={blue},
  pdfcreator={LaTeX via pandoc}}

\title{RP8 --- The source code of Revision\_Deflation: what computes
what, and how to compile and run it}
\usepackage{etoolbox}
\makeatletter
\providecommand{\subtitle}[1]{% add subtitle to \maketitle
  \apptocmd{\@title}{\par {\large #1 \par}}{}{}
}
\makeatother
\subtitle{Provenance required by Hard Rule 6 (PROJECT spacetimes,
Revision\_Deflation)}
\author{Prepared for Patrick Lee Nash, Ph.D.}
\date{November 13, 2026}

\begin{document}
\maketitle

\hypertarget{what-this-document-is-for}{%
\section{1. What this document is for}\label{what-this-document-is-for}}

Hard Rule 6 demands commented source code for every equation, result,
figure and prediction, instructions for using it, a provenance for every
computed value, and that the code be executed. This document is the map
of the code of the revision: which program computes which file, which
file feeds which figure and which table, how the programs are tested,
and how to compile and run everything. It follows the provenance P8 of
the main project, which the reader may consult for the general rules;
they are the same here:

\begin{itemize}
\tightlist
\item
  \textbf{Rust for every computation.} All physics of the revision is
  computed by one Rust program, \texttt{deflation} (folder
  \texttt{code/deflation-rs}: 24 source files, 10324 lines, 49 unit
  tests). It uses the program of the main project as a library and no
  other library. Four of its source files are copies of modules of the
  main project with a listed set of changes; they are written by a
  script and not edited by hand (Sec. 3.1).
\item
  \textbf{Python for an independent verification, and for tables,
  figures and documents.} The central results are verified a second time
  by a Python script that shares no code with the Rust program and uses
  other methods (computer algebra instead of rational arithmetic; a
  general-purpose integrator and library Bessel functions instead of the
  integrator of the program).
\item
  \textbf{No number is typed by hand.} Every table of a document is
  written by a script from a data file that the program wrote; every
  check quoted in a document is reproduced verbatim in an appendix.
\end{itemize}

\textbf{Where the code was executed.} On a Linux computer (x86-64) with
rustc 1.97.0, Python 3.13 (numpy 2.5.3, scipy 1.18.1, sympy 1.14.0,
mpmath 1.3.0, matplotlib 3.11.2), pandoc and XeLaTeX. As for the main
project, the assistant could not run the Rust program on the author's
Mac, where no Rust compiler is reachable from its sandbox; the program
has no dependency outside the repository and should build and run on
macOS unchanged.

\hypertarget{how-to-compile-and-run-everything-on-a-mac}{%
\section{2. How to compile and run everything on a
Mac}\label{how-to-compile-and-run-everything-on-a-mac}}

Each block can be pasted into the Terminal as it is. A line that begins
with \texttt{\#} is an explanation and does nothing. The Rust compiler
and the Python packages are installed as described in P8 of the main
project, Sec. 2, Steps 1 and 2; if the main project has been run on this
Mac, nothing needs to be installed. The independent check needs two
Python packages more than the main project:

\begin{Shaded}
\begin{Highlighting}[]
\CommentTok{\# once: the two additional Python packages (computer algebra; precise special functions)}
\ExtensionTok{python3} \AttributeTok{{-}m}\NormalTok{ pip install }\AttributeTok{{-}{-}user}\NormalTok{ sympy mpmath}
\end{Highlighting}
\end{Shaded}

\textbf{Step 1: go to the folder of the revision and build the program.}

\begin{Shaded}
\begin{Highlighting}[]
\CommentTok{\# every command of the revision is run from this folder}
\BuiltInTok{cd} \StringTok{"/Users/nsh/Developer/github/spacetimes/Revision\_Deflation"}
\CommentTok{\# are the four copied modules what vendor\_eb.py writes from the sources of the main project?}
\CommentTok{\# (must end with "the four copied modules are up to date";}
\CommentTok{\#  without {-}{-}check the script writes them again)}
\ExtensionTok{python3}\NormalTok{ code/scripts/vendor\_eb.py }\AttributeTok{{-}{-}check}
\CommentTok{\# compile the program (the first time also that of the main project; about a minute)}
\ExtensionTok{cargo}\NormalTok{ build }\AttributeTok{{-}{-}release} \AttributeTok{{-}{-}manifest{-}path}\NormalTok{ code/deflation{-}rs/Cargo.toml}
\CommentTok{\# run the unit tests; the last lines must say "test result: ok. 49 passed"}
\ExtensionTok{cargo}\NormalTok{ test }\AttributeTok{{-}{-}release} \AttributeTok{{-}{-}manifest{-}path}\NormalTok{ code/deflation{-}rs/Cargo.toml}
\end{Highlighting}
\end{Shaded}

\textbf{Step 2: run the fourteen commands.} Each prints its checks, one
per line, beginning with \texttt{{[}PASS{]}} or \texttt{{[}FAIL{]}}, and
ends with a line such as \texttt{matrices:\ ...\ checks\ passed}. The
first nine take about three and a half minutes together. The last five
are the Einstein--Boltzmann calculations of the supplement RS5 and take
about an hour and ten minutes on two processor cores, most of it the
commands \texttt{families} and \texttt{superfluid}. Together the
fourteen rewrite every file of the folder \texttt{data}.

\begin{Shaded}
\begin{Highlighting}[]
\CommentTok{\# geometry, Dirac matrices, spin connection, Dirac operators           (less than a second)}
\ExtensionTok{./code/deflation{-}rs/target/release/deflation}\NormalTok{ matrices}
\CommentTok{\# solutions of the Dirac equations: creation and extinction           (about 12 seconds)}
\ExtensionTok{./code/deflation{-}rs/target/release/deflation}\NormalTok{ modes}
\CommentTok{\# gap, bound pair, matter, vacuum, temperature along the clock        (about 1.5 minutes)}
\ExtensionTok{./code/deflation{-}rs/target/release/deflation}\NormalTok{ history}
\CommentTok{\# can the paired state follow the deflation?                          (about 16 seconds)}
\ExtensionTok{./code/deflation{-}rs/target/release/deflation}\NormalTok{ quench}
\CommentTok{\# Hubble constant, bounds on the rate, the exponent under H2          (about a minute)}
\ExtensionTok{./code/deflation{-}rs/target/release/deflation}\NormalTok{ cosmology}
\CommentTok{\# surviving symmetries, vortex zero modes, exponents of the couplings (less than a second)}
\ExtensionTok{./code/deflation{-}rs/target/release/deflation}\NormalTok{ symmetry}
\CommentTok{\# exchange amplitudes, modes of the scalar                            (about a second)}
\ExtensionTok{./code/deflation{-}rs/target/release/deflation}\NormalTok{ channels}
\CommentTok{\# dark matter, growth of structure, drifts                            (about 10 seconds)}
\ExtensionTok{./code/deflation{-}rs/target/release/deflation}\NormalTok{ matter}
\CommentTok{\# the paired state from the effective action                          (about 10 seconds)}
\ExtensionTok{./code/deflation{-}rs/target/release/deflation}\NormalTok{ paired}
\CommentTok{\# recombination and the sound horizon on the deflating background     (about 5 seconds)}
\ExtensionTok{./code/deflation{-}rs/target/release/deflation}\NormalTok{ background}
\CommentTok{\# spectra of the microwave background at a fixed angle; closures      (about a minute)}
\ExtensionTok{./code/deflation{-}rs/target/release/deflation}\NormalTok{ spectrum}
\CommentTok{\# mock data against the rate of deflation: Fisher matrices and fits   (about 6 minutes)}
\ExtensionTok{./code/deflation{-}rs/target/release/deflation}\NormalTok{ forecast}
\CommentTok{\# the dark{-}radiation families of the main project, with deflation     (about 40 minutes)}
\ExtensionTok{./code/deflation{-}rs/target/release/deflation}\NormalTok{ families}
\CommentTok{\# superfluid dark matter: suppression of small{-}scale structure        (about 20 minutes)}
\ExtensionTok{./code/deflation{-}rs/target/release/deflation}\NormalTok{ superfluid}
\end{Highlighting}
\end{Shaded}

(\texttt{deflation\ families\ quick} runs two of the seven families at
two densities for one data set in three minutes. It is a test of the
command: it writes its files to a temporary folder, which it names, and
leaves the folder \texttt{data} as it is.)

\textbf{Step 3: the independent verification} (about 20 seconds; must
end with ``43 of 43 checks passed''), and the test that it can fail
(about four minutes; must end with ``16 of 16 deliberate errors
detected'').

\begin{Shaded}
\begin{Highlighting}[]
\ExtensionTok{python3}\NormalTok{ code/scripts/verify\_deflating\_independent.py}
\ExtensionTok{python3}\NormalTok{ code/scripts/mutation\_test\_independent.py}
\end{Highlighting}
\end{Shaded}

\textbf{Step 4: symbols, tables, figures, documents.}

\begin{Shaded}
\begin{Highlighting}[]
\CommentTok{\# the registry of symbols (provenance RP1)}
\ExtensionTok{python3}\NormalTok{ code/scripts/write\_symbol\_registry.py}
\CommentTok{\# are the numbers quoted in the running text of the documents numbers that the programs stated?}
\ExtensionTok{python3}\NormalTok{ code/scripts/check\_numbers.py}
\CommentTok{\# the tables of the documents, from the data files and from the reports of the checks}
\ExtensionTok{python3}\NormalTok{ code/scripts/write\_doc\_tables.py}
\CommentTok{\# the 19 figures (written to the folder figures as .pdf and .png)}
\ExtensionTok{python3}\NormalTok{ code/scripts/make\_figures.py}
\CommentTok{\# the documents: Markdown from the sources, then LaTeX and PDF (needs pandoc and XeLaTeX)}
\ExtensionTok{python3}\NormalTok{ code/scripts/assemble\_docs.py}
\FunctionTok{bash}\NormalTok{ code/scripts/build\_docs.sh}
\CommentTok{\# the two manuscripts (needs pdflatex with REVTeX 4.2)}
\ExtensionTok{python3}\NormalTok{ code/scripts/build\_manuscripts.py}
\CommentTok{\# the six checks of the symbols (they read the registry, the data files and the documents)}
\ExtensionTok{python3}\NormalTok{ code/scripts/check\_symbols.py}
\end{Highlighting}
\end{Shaded}

Two of these scripts read documents and write reports that documents
print: the report of \texttt{check\_numbers.py} is quoted in RP8 and in
the README, and the report of \texttt{check\_symbols.py} is printed in
RP1. After a change of a document the commands of Step 4 are therefore
run twice, and the second run must leave every data file, report, table,
figure, Markdown file and LaTeX file as it was. (The PDF files of the
documents other than the manuscripts carry the time of their compilation
and differ in that alone.) In the repository as it stands a single run
leaves these files as they are.

If a command of Step 2 prints a line that begins with
\texttt{{[}FAIL{]}}, or if the numbers in a document change when Steps 2
to 4 are repeated, something is wrong and should be reported; on another
processor the last digits of floating-point results may differ, the
exact checks may not.

\hypertarget{the-rust-program-file-by-file}{%
\section{3. The Rust program, file by
file}\label{the-rust-program-file-by-file}}

All files are in \texttt{code/deflation-rs/src}. Every file begins with
a description of the physics that it computes, with the formulas; every
public function is documented.

\begin{longtable}[]{@{}
  >{\raggedright\arraybackslash}p{(\columnwidth - 4\tabcolsep) * \real{0.2105}}
  >{\raggedleft\arraybackslash}p{(\columnwidth - 4\tabcolsep) * \real{0.2105}}
  >{\raggedright\arraybackslash}p{(\columnwidth - 4\tabcolsep) * \real{0.5789}}@{}}
\toprule\noalign{}
\begin{minipage}[b]{\linewidth}\raggedright
file
\end{minipage} & \begin{minipage}[b]{\linewidth}\raggedleft
lines
\end{minipage} & \begin{minipage}[b]{\linewidth}\raggedright
content
\end{minipage} \\
\midrule\noalign{}
\endhead
\bottomrule\noalign{}
\endlastfoot
\texttt{rat.rs} & 223 & exact rational numbers with overflow checks; the
trait that lets the geometry be computed in rational or in
floating-point arithmetic \\
\texttt{qmat.rs} & 201 & exact matrices with complex-rational entries \\
\texttt{metric.rs} & 714 & the deflating metric: Christoffel symbols,
Riemann, Ricci and Einstein tensors from their definitions and in closed
form; Killing vectors; Lie derivative \\
\texttt{geometry.rs} & 443 & the exact geometry checks \texttt{G1} to
\texttt{G7}, \texttt{G10} (Lovelock tensors) and the geodesics
\texttt{G8}, \texttt{G9} \\
\texttt{curved.rs} & 530 & coordinate matrices, spin connection, Dirac
operators; checks \texttt{C1} to \texttt{C8}, \texttt{D1} to
\texttt{D7} \\
\texttt{modes.rs} & 766 & integration of the mode equations
(commutator-free Magnus method), creation of particles, bounded
solutions, modified Bessel function of complex order \\
\texttt{scalar.rs} & 295 & modes of a scalar field: exponents, static
source, creation, bounded solutions \\
\texttt{history.rs} & 274 & the couplings along the clock; thresholds
and events \\
\texttt{quench.rs} & 280 & time-dependent mean-field dynamics of the
shell model \\
\texttt{cosmo\_defl.rs} & 512 & expansion with deflation, sound horizon,
age, bounds on the rate, growth of pressureless matter for the
closures \\
\texttt{spectra\_defl.rs} & 611 & assumption A3 and its closures; the
Hubble constant at a fixed angle with the full thermal history;
potentials and gaps of single modes; derivatives of the spectra \\
\texttt{eb\_cosmo.rs}, \texttt{eb\_boltzmann.rs},
\texttt{eb\_spectra.rs}, \texttt{eb\_fit.rs} & 1891 & the
Einstein-Boltzmann program of the main project on the deflating
spacetime: copies of four of its modules with a listed set of changes,
written by \texttt{code/scripts/vendor\_eb.py} \\
\texttt{symmetry\_defl.rs} & 217 & test of the 28 rotations and boosts;
algebra of the vortex zero modes \\
\texttt{cmd\_symmetry.rs}, \texttt{cmd\_channels.rs},
\texttt{cmd\_matter.rs}, \texttt{cmd\_paired.rs} & 1017 & the commands
\texttt{symmetry}, \texttt{channels}, \texttt{matter},
\texttt{paired} \\
\texttt{cmd\_spectra.rs} & 1425 & the commands \texttt{background},
\texttt{spectrum}, \texttt{forecast}, \texttt{families},
\texttt{superfluid} \\
\texttt{main.rs} & 808 & the commands \texttt{matrices}, \texttt{modes},
\texttt{history}, \texttt{quench}, \texttt{cosmology} \\
\texttt{io.rs}, \texttt{lib.rs} & 117 & writing of data files and
reports; list of modules \\
\end{longtable}

(The numbers of lines and of unit tests in this document are counted by
the script that writes the tables; they are not typed.)

\textbf{What is taken from the program of the main project} (crate
\texttt{spacetimes}, used as a library and not changed): the Dirac
matrices and their checks, the routine for the internal symmetry, the
exchange amplitudes, the gap equation, the bound pair, the effective
action, the routines for matter, temperature and vacuum, the background
cosmology without deflation, linear algebra, and from its
Einstein--Boltzmann part the parameters, the dark fluids, the superfluid
dark matter, the Bessel functions, the mock data, the likelihood and the
Fisher matrix. The revision calls these routines with the couplings of
each moment; it does not copy them.

\hypertarget{the-four-copied-modules}{%
\subsection{3.1 The four copied modules}\label{the-four-copied-modules}}

The Einstein--Boltzmann program of the main project has Newton's
constant built into its units and the ordinary equation of the expansion
built into its routines. On the deflating spacetimes both change (RS3,
Sec. 1; RS5, Sec. 2). The main project is paused and its program is not
to be changed, so the revision cannot add a parameter to it; and a
library routine cannot be altered from outside. The four modules that
must change are therefore \emph{copied}:

\begin{longtable}[]{@{}
  >{\raggedright\arraybackslash}p{(\columnwidth - 4\tabcolsep) * \real{0.1905}}
  >{\raggedright\arraybackslash}p{(\columnwidth - 4\tabcolsep) * \real{0.1905}}
  >{\raggedright\arraybackslash}p{(\columnwidth - 4\tabcolsep) * \real{0.6190}}@{}}
\toprule\noalign{}
\begin{minipage}[b]{\linewidth}\raggedright
file of the revision
\end{minipage} & \begin{minipage}[b]{\linewidth}\raggedright
copied from (main project, \texttt{code/spacetimes-rs/src})
\end{minipage} & \begin{minipage}[b]{\linewidth}\raggedright
what is changed
\end{minipage} \\
\midrule\noalign{}
\endhead
\bottomrule\noalign{}
\endlastfoot
\texttt{eb\_cosmo.rs} & \texttt{cosmo.rs}: the structures
\texttt{Background} and \texttt{Thermo} & the equation of the expansion
(RS3, Eq. (1.2)); the clock of Newton's constant; the helium fraction \\
\texttt{eb\_boltzmann.rs} & \texttt{boltzmann.rs} & Newton's constant of
the moment in every source of the metric; \(H_e\) in the energy
constraint; the closure ``Pressure''; the gaps; the sources of the
temperature \\
\texttt{eb\_spectra.rs} & \texttt{spectra.rs}: the functions
\texttt{compute} and \texttt{matter\_power} & they take the deflation as
an argument \\
\texttt{eb\_fit.rs} & \texttt{forecast.rs}: the functions
\texttt{minimize\_linearized} and \texttt{fit} & they take the deflation
as an argument \\
\end{longtable}

So that a copy cannot differ from its original in any place that is not
listed, the copies are \textbf{not edited by hand}. The script
\texttt{code/scripts/vendor\_eb.py} reads the four original files and
applies a list of replacements, each of which must match the original
text exactly once; otherwise it stops. The head of the script states
every change in words, and the head of each written file repeats the
list. If the main project changes one of the four files,
\texttt{vendor\_eb.py\ -\/-check} reports that the copies are no longer
current, and running the script without \texttt{-\/-check} either
reproduces the copies from the new originals or stops at the replacement
that no longer matches.

Three tests guard the copies (file \texttt{spectra\_defl.rs}; checks
\texttt{def-E1}, \texttt{def-B1}, \texttt{def-F1}, \texttt{def-U1}):

\begin{enumerate}
\def\labelenumi{\arabic{enumi}.}
\tightlist
\item
  \emph{Without deflation the copies are the program of the main
  project.} The spectra TT, TE, EE for \(2\le\ell\le2500\), the matter
  power spectrum, \(\sigma_8\), the angle and the sound horizon agree
  \textbf{bit for bit} with the library routines, for the fiducial
  model, for superfluid dark matter and for a stepped dark fluid, and
  for both closures.
\item
  \emph{With a rate too small to matter the deflating code paths give
  the same numbers.} A rate of \(10^{-30}\) (in units of 100 km/s/Mpc)
  makes the program execute every branch of the deflating model;
  potentials, sources and spectra are again identical to the last bit.
\item
  \emph{With a real rate each closure is a consistent theory.} The
  quantity that the source of the temperature uses equals the derivative
  of the potentials along the solution, and the two identities between
  the gaps (RS5, Sec. 3.1) hold along the solution to the accuracy of a
  numerical derivative.
\end{enumerate}

\hypertarget{the-checks}{%
\section{4. The checks}\label{the-checks}}

\begin{longtable}[]{@{}
  >{\raggedright\arraybackslash}p{(\columnwidth - 4\tabcolsep) * \real{0.3333}}
  >{\raggedright\arraybackslash}p{(\columnwidth - 4\tabcolsep) * \real{0.3333}}
  >{\raggedleft\arraybackslash}p{(\columnwidth - 4\tabcolsep) * \real{0.3333}}@{}}
\toprule\noalign{}
\begin{minipage}[b]{\linewidth}\raggedright
command
\end{minipage} & \begin{minipage}[b]{\linewidth}\raggedright
file of checks
\end{minipage} & \begin{minipage}[b]{\linewidth}\raggedleft
number
\end{minipage} \\
\midrule\noalign{}
\endhead
\bottomrule\noalign{}
\endlastfoot
\texttt{matrices} & \texttt{data/flat\_matrix\_checks.txt} (the checks
of the main project, again) & 106 of 106 \\
\texttt{matrices} & \texttt{data/deflating\_checks.txt} & 47 of 47 \\
\texttt{modes} & \texttt{data/mode\_checks.txt} & 7 of 7 \\
\texttt{history} & \texttt{data/history\_checks.txt} & 8 of 8 \\
\texttt{quench} & \texttt{data/quench\_checks.txt} & 5 of 5 \\
\texttt{cosmology} & \texttt{data/cosmology\_checks.txt} & 6 of 6 \\
\texttt{symmetry} & \texttt{data/symmetry\_checks.txt} & 18 of 18 \\
\texttt{channels} & \texttt{data/channel\_checks.txt} & 7 of 7 \\
\texttt{matter} & \texttt{data/matter\_checks.txt} & 6 of 6 \\
\texttt{paired} & \texttt{data/paired\_checks.txt} & 7 of 7 \\
\texttt{background} & \texttt{data/background\_checks.txt} & 2 of 2 \\
\texttt{spectrum} & \texttt{data/spectrum\_checks.txt} & 8 of 8 \\
\texttt{forecast} & \texttt{data/forecast\_checks.txt} & 5 of 5 \\
\texttt{families} & \texttt{data/families\_checks.txt} & 4 of 4 \\
\texttt{superfluid} & \texttt{data/superfluid\_checks.txt} & 2 of 2 \\
Python script & \texttt{data/deflating\_checks\_independent.txt} & 43 of
43 \\
symbols & \texttt{data/symbol\_checks\_deflation.txt} & 6 of 6 \\
numbers of the running text &
\texttt{data/number\_checks\_deflation.txt} & 974 numbers in 29
documents, 0 not traced \\
\end{longtable}

Each file of checks exists also as \texttt{.csv}. A check marked
``exact'' is an identity verified in exact rational arithmetic; a check
marked ``numerical'' states its accuracy. The identifiers begin with
\texttt{44d-} (deflating 4+4), \texttt{71d-} (deflating 7+1) or
\texttt{def-} (both); those of the Python script with \texttt{I}; those
of the symbols with \texttt{RS}.

\emph{Was the independent script itself tested?} Yes: twenty deliberate
errors were introduced into it, one at a time (a wrong sign in a
Christoffel symbol, a wrong factor in the Einstein tensor, the wrong
order of a Bessel function, and so on). The script detected nineteen.
The twentieth, a change of the term \(\mathrm{gprod}\) in the equation
of the expansion, cannot be detected by the comparison of Hubble
constants, because that term is absorbed by the cosmological term (RS3,
Sec. 2); it is tested by the symbolic check of the Einstein tensor
instead. The checks of the completeness of the Killing vectors and of
the Lanczos tensor were added afterwards; the second is validated by the
agreement of two independent formulas (the program uses the formula for
curvatures that act on pairs of directions, the script the standard
formula with four terms).

The checks of the perturbation equations, of the closures, of the growth
of structure and of recombination, which were added for the supplement
RS5, were tested in the same way, and for them the test is itself a
script: \texttt{mutation\_test\_independent.py} introduces sixteen
deliberate errors, one at a time, into a copy of that part of the
independent script (a wrong factor in an Einstein equation, in an
identity between the gaps, in a growth equation, in the response of
recombination and of the helium fraction to Newton's constant, and so
on), runs the copy, and records which check fails. Its report:

\begin{Shaded}
\begin{Highlighting}[]
\NormalTok{PROJECT spacetimes, Revision\_Deflation: mutation test of the independent verification of the cosmological perturbations}
\NormalTok{generated by: code/scripts/mutation\_test\_independent.py}

\NormalTok{the unchanged script: 11 checks (Idef{-}C1, Idef{-}C2, I44d{-}P1, I71d{-}P1, Idef{-}P2, Idef{-}P8, Idef{-}P3, Idef{-}P4, Idef{-}P5, Idef{-}P6, Idef{-}P7), all pass}

\NormalTok{ 1. energy constraint with He = Hc {-} a gsum instead of Hc {-} a gsum/2: detected by I44d{-}P1, I71d{-}P1}
\NormalTok{ 2. trace equation with He h\textquotesingle{} instead of 2 He h\textquotesingle{}: detected by I44d{-}P1, I71d{-}P1}
\NormalTok{ 3. closed form of B with (a gsum)\^{}2/2 instead of (a gsum)\^{}2/4: detected by Idef{-}P2}
\NormalTok{ 4. pressure of the agent with gprod instead of 2 gprod: detected by Idef{-}P8}
\NormalTok{ 5. first identity between the gaps with (D {-} DP) instead of (D {-} 3 DP): detected by Idef{-}P3}
\NormalTok{ 6. second identity between the gaps with He instead of 2 He: detected by Idef{-}P3}
\NormalTok{ 7. growth equation of the closure "Momentum" with the source not reduced: detected by Idef{-}P4}
\NormalTok{ 8. growth equation of the closure "Pressure" with the friction Hc instead of He: detected by Idef{-}P4}
\NormalTok{ 9. numerical growth, closure "Pressure": friction reduced by gsum instead of gsum/2: detected by Idef{-}P5, Idef{-}P7}
\NormalTok{10. numerical growth, closure "Momentum": source not reduced: detected by Idef{-}P5, Idef{-}P7}
\NormalTok{11. system of the closure "Pressure" with a gsum instead of a gsum/2: detected by Idef{-}P6}
\NormalTok{12. system of the closure "Shear" with Hc instead of He: detected by Idef{-}P6}
\NormalTok{13. recombination on a background in which Newton\textquotesingle{}s constant does not change: detected by Idef{-}C2}
\NormalTok{14. response of the helium fraction to Newton\textquotesingle{}s constant doubled: detected by Idef{-}C2}
\NormalTok{15. Peebles\textquotesingle{} factor with 4 pi instead of 8 pi (an error of atomic physics): detected by Idef{-}C2}
\NormalTok{16. equation of the expansion with the term {-}gsum H doubled (integration at fixed last scattering): detected by Idef{-}C1}

\NormalTok{16 of 16 deliberate errors detected}
\end{Highlighting}
\end{Shaded}

\emph{Are the numbers in the text the numbers of the programs?} The
tables of the documents are written by a script. The running text quotes
numbers by hand, and a quoted number can be mistyped or can outlive a
change of the program. The script \texttt{check\_numbers.py} extracts
every number from the running text of the sources of all documents and
of the two manuscripts (numbers with a decimal point, numbers in
scientific notation, and integers from 20 on) and looks it up among the
numbers that the programs \emph{stated}: in the statements of the
checks, in the generated tables, and in the descriptions at the head of
the data files. A number of the text counts as traced if

\begin{itemize}
\item
  \begin{enumerate}
  \def\labelenumi{(\alph{enumi})}
  \tightlist
  \item
    a stated number equals it; or
  \end{enumerate}
\item
  \begin{enumerate}
  \def\labelenumi{(\alph{enumi})}
  \setcounter{enumi}{17}
  \tightlist
  \item
    a stated number rounds to it at the precision typed; or
  \end{enumerate}
\item
  \begin{enumerate}
  \def\labelenumi{(\alph{enumi})}
  \setcounter{enumi}{1}
  \tightlist
  \item
    it is listed in the script, with the reason, as a number that the
    programs do not state: a value from the literature, a definition, a
    count that follows from a formula of the text, a parameter of a run,
    or (in the error log RP9 only) a superseded value that the log
    quotes; or
  \end{enumerate}
\item
  \begin{enumerate}
  \def\labelenumi{(\alph{enumi})}
  \setcounter{enumi}{2}
  \tightlist
  \item
    it is typed with one significant digit in scientific notation, as a
    bound or a power of ten, and a stated number lies below it within
    one decade.
  \end{enumerate}
\end{itemize}

The numbers inside the data files are deliberately not used: the files
hold more than a hundred thousand numbers, and almost any number is the
rounded value of one of them by coincidence. The matches of the kinds
(r) and (c) are weaker than those of kind (a): among several thousand
stated numbers a short number can have a rounded partner by coincidence.
Each number of these two kinds must therefore be \emph{judged} by a
reader, who reads the sentence and the output of the programs and enters
into the script the stated number that the text means. The report
\texttt{data/number\_checks\_deflation.txt} lists these numbers one by
one, with the stated numbers that match and with the judgement, and the
script fails if one of them has not been judged. Result: 974 numbers in
29 documents, 0 not traced; different numbers of the kinds (r) and (c),
all judged: 10. The check is a necessary condition, not a sufficient
one: it cannot tell whether the right number is quoted in the right
place, and a match can be a coincidence. It found the errors R-E28 and
R-E29 of the error log RP9. The stale value of R-E27 passed it, because
a number of a generated table rounds to it by coincidence; that error
was found by reading the program.

\hypertarget{the-scripts}{%
\section{5. The scripts}\label{the-scripts}}

All in \texttt{code/scripts}.

\begin{longtable}[]{@{}
  >{\raggedright\arraybackslash}p{(\columnwidth - 2\tabcolsep) * \real{0.2353}}
  >{\raggedright\arraybackslash}p{(\columnwidth - 2\tabcolsep) * \real{0.7647}}@{}}
\toprule\noalign{}
\begin{minipage}[b]{\linewidth}\raggedright
script
\end{minipage} & \begin{minipage}[b]{\linewidth}\raggedright
what it does
\end{minipage} \\
\midrule\noalign{}
\endhead
\bottomrule\noalign{}
\endlastfoot
\texttt{vendor\_eb.py} & writes the four copied modules of the
Einstein--Boltzmann program from the sources of the main project (Sec.
3.1); with \texttt{-\/-check}, tests that they are current \\
\texttt{verify\_deflating\_independent.py} & the independent
verification (Sec. 4) \\
\texttt{mutation\_test\_independent.py} & introduces sixteen deliberate
errors into the cosmological part of the independent verification, one
at a time, and records which check detects each (Sec. 4) \\
\texttt{write\_symbol\_registry.py}, \texttt{check\_symbols.py} & the
registry of the symbols and its six checks (RP1) \\
\texttt{write\_doc\_tables.py} & writes every table of the documents
from the data files
(\texttt{provenance/generated/tables\_deflation.md}) \\
\texttt{make\_figures.py} & draws the 19 figures from the data files \\
\texttt{assemble\_docs.py} & inserts tables, registry and program output
into the document sources (\texttt{docs\_src/*.src.md}) \\
\texttt{build\_docs.sh}, \texttt{doc\_header.tex},
\texttt{check\_layout.py} & Markdown to LaTeX to PDF, and a check that
no line, table cell or formula is wider than its box \\
\texttt{build\_manuscripts.py}, \texttt{tex2md.py} & the two manuscripts
from one source (\texttt{docs\_src/manuscript/*.tex}); their Markdown
versions \\
\texttt{check\_numbers.py} & looks up every number of the running text
of the documents in the output of the programs (Sec. 4) \\
\end{longtable}

\hypertarget{which-data-file-feeds-which-figure}{%
\section{6. Which data file feeds which
figure}\label{which-data-file-feeds-which-figure}}

\begin{longtable}[]{@{}
  >{\raggedright\arraybackslash}p{(\columnwidth - 4\tabcolsep) * \real{0.2222}}
  >{\raggedright\arraybackslash}p{(\columnwidth - 4\tabcolsep) * \real{0.5556}}
  >{\raggedright\arraybackslash}p{(\columnwidth - 4\tabcolsep) * \real{0.2222}}@{}}
\toprule\noalign{}
\begin{minipage}[b]{\linewidth}\raggedright
figure
\end{minipage} & \begin{minipage}[b]{\linewidth}\raggedright
data files
\end{minipage} & \begin{minipage}[b]{\linewidth}\raggedright
in
\end{minipage} \\
\midrule\noalign{}
\endhead
\bottomrule\noalign{}
\endlastfoot
\texttt{fig\_geodesics} & \texttt{geodesic\_1671d.csv},
\texttt{geodesic\_1644d.csv} & DMD, M \\
\texttt{fig\_isometries} & \texttt{isometries.csv} & DMD, RS2, M \\
\texttt{fig\_creation} & \texttt{creation\_single.csv},
\texttt{creation\_history.csv}, \texttt{creation\_anisotropic.csv},
\texttt{boson\_creation.csv} & DMD, M (7+1) \\
\texttt{fig\_extinction} & \texttt{extinction\_profile.csv},
\texttt{boson\_bounded.csv} & DMD, M (4+4) \\
\texttt{fig\_clock} & \texttt{coupling\_exponents.csv},
\texttt{events.csv} & RS1, M \\
\texttt{fig\_scalar} & \texttt{scalar\_zero\_modes.csv},
\texttt{static\_source.csv} & RS1, M \\
\texttt{fig\_kernels} & \texttt{pair\_channel\_kernel.csv} & RS1, M \\
\texttt{fig\_gap\_history} & \texttt{gap\_history.csv} & RS1, M \\
\texttt{fig\_binding} & \texttt{dimer\_history.csv},
\texttt{dm\_resonance.csv} & RS1, RS3, M \\
\texttt{fig\_quench} & \texttt{quench\_history.csv},
\texttt{quench\_summary.csv} & RS1, M \\
\texttt{fig\_cjt\_history} & \texttt{cjt\_history.csv} & RS4, M \\
\texttt{fig\_matter\_history} & \texttt{medium\_history.csv},
\texttt{matter\_saturation.csv}, \texttt{vacuum\_history.csv},
\texttt{thermal\_history.csv} & RS4, M \\
\texttt{fig\_hubble} & \texttt{hubble\_deflation.csv},
\texttt{hubble\_mass\_loss.csv}, \texttt{growth\_deflation.csv} & RS3,
M \\
\texttt{fig\_bounds} & \texttt{deflation\_bounds.csv},
\texttt{deflation\_fisher.csv} & RS3, RS2, M \\
\texttt{fig\_spectra} & \texttt{cmb\_deflation\_cl.csv},
\texttt{cmb\_deflation\_summary.csv} & RS5, M \\
\texttt{fig\_closure} & \texttt{cmb\_deflation\_closure.csv},
\texttt{cmb\_deflation\_gaps.csv}, \texttt{cmb\_deflation\_growth.csv},
\texttt{growth\_deflation.csv} & RS5, M \\
\texttt{fig\_forecast} & \texttt{deflation\_fisher.csv},
\texttt{deflation\_profile.csv}, \texttt{deflation\_residuals.csv},
\texttt{hubble\_deflation.csv} & RS5, M \\
\texttt{fig\_families} & \texttt{dr\_deflation\_profiles.csv},
\texttt{dr\_deflation\_joint.csv}, and \texttt{dr\_profiles.csv} of the
main project & RS5 \\
\texttt{fig\_sfdm} & \texttt{sfdm\_deflation\_transfer.csv},
\texttt{sfdm\_deflation\_summary.csv} & RS5 \\
\end{longtable}

(M: the two manuscripts.) Every data file begins with lines that state
which program and which command wrote it and what its columns mean.

\end{document}
